% Margin glosses for acronyms and jargon (\firstterm{key} on first use).
\makeatletter
\newcommand{\declaretermsnip}[2]{\expandafter\gdef\csname termsnip@#1\endcsname{#2}}
\newcommand{\firstterm}[1]{%
  \expandafter\ifx\csname termused@#1\endcsname\relax
    \expandafter\gdef\csname termused@#1\endcsname{}%
    \ifcsname termsnip@#1\endcsname
      \sidenote{\csname termsnip@#1\endcsname}%
    \fi
  \fi
}
\makeatother

\declaretermsnip{FEC}{Forward error correction: parity symbols let the receiver fix
  bit errors without retransmit. Datacom links specify a \emph{pre-FEC} BER the code
  must correct.}

\declaretermsnip{KP4}{IEEE 802.3 Clause~91 FEC for PAM4: Reed--Solomon RS(544,514) on
  10-bit symbols, 5.5\% overhead. Pre-FEC BER target $2.4\times10^{-4}$.}

\declaretermsnip{RS544514}{Reed--Solomon block code: 544 symbols out, 514 payload, 30
  parity. Corrects up to 15 symbol errors per codeword.}

\declaretermsnip{preFEC}{Pre-FEC BER: measured before the FEC decoder. Standards and
  validation sweeps use this number, not post-FEC.}

\declaretermsnip{postFEC}{Post-FEC BER: after FEC correction. Target is effectively
  error-free ($<10^{-12}$ class).}

\declaretermsnip{MLC}{Multi-level coding: coded modulation paired with PAM6/PAM8. OIF
  448G workshop models often add a strong outer RS code; $\sim$12\% overhead vs.\ KP4
  at 5.5\%.}

\declaretermsnip{HDFEC}{Hard-decision FEC: high-overhead codes used in optical demos and
  long-reach products (10--15\% OH), not the host KP4 budget.}

\declaretermsnip{OH}{Coding overhead: fraction of line rate spent on parity symbols
  (KP4 $\approx$5.5\%).}

\declaretermsnip{TDECQ}{Transmitter and dispersion eye closure quaternary (TDECQ):
  headline PAM4 transmitter-quality metric after a reference receiver and bounded
  FFE, including a test fiber. Lower is better; LPO MSA 100G-DR caps near 3.4~dB.
  Replaced simple eye-mask tests. Stressed counterpart: SECQ.}

\declaretermsnip{TECQ}{Transmitter eye closure quaternary (TECQ): same measurement as
  TDECQ without the test fiber. Used with TDECQ in LPO MSA OMA tables.}

\declaretermsnip{SECQ}{Stressed eye closure quaternary: receiver-side stressed test,
  analog of TDECQ. Applies calibrated optical stress before BER test.}

\declaretermsnip{EAM}{Electro-absorption modulator: voltage-controlled absorption on InP;
  low chirp vs DML; paired with DFB in EML.}

\declaretermsnip{APD}{Avalanche photodiode: internal gain improves sensitivity; excess
  noise factor and bias complexity vs PIN.}

\declaretermsnip{TIA}{Transimpedance amplifier: converts PD current to voltage; input
  capacitance sets input-referred noise.}

\declaretermsnip{PIN}{PIN photodiode: no internal gain; Ge-on-Si mainstream at 100--224G.}

\declaretermsnip{gearbox}{Rate/format converter between electrical and optical lanes;
  common in 448G transition; incompatible with strict LPO.}

\declaretermsnip{CMIS}{Common Management Interface Specification: module management,
  monitors, link training at 224G/448G.}

\declaretermsnip{RMA}{Return merchandise authorization: field-failed unit returned for
  FA. Distinct failure codes keep fleet FIT honest.}

\declaretermsnip{ATP}{Acceptance test plan: measurable ship tests that prove the
  requirements contract with a supplier (often mapped to GR-468 stresses).}

\declaretermsnip{SPC}{Statistical process control: control charts on key parameters
  ($I_\mathrm{th}$, SMSR, power, burn-in fallout) to catch drift qual lots missed.}

\declaretermsnip{eightD}{Eight disciplines (8D) corrective action: contain, root-cause,
  correct, and prevent with the supplier; close only with verification evidence.}

\declaretermsnip{DPA}{Destructive physical analysis: cross-section / EDX on failed
  units to confirm facet, solder, or FAU failure modes.}

\declaretermsnip{LOS}{Loss of signal: receiver (or host) flag that optical input is
  below the detect threshold.}

\declaretermsnip{LOL}{Loss of lock: CDR or SerDes unlocked; electrical path not
  recovered even if light is present.}

\declaretermsnip{FAIR}{First-article inspection report: re-qualify after tooling, site,
  epi, or CMIS firmware change before open volume.}

\declaretermsnip{DML}{Directly modulated laser: modulate bias current; high chirp, low cost.}

\declaretermsnip{VCSEL}{Vertical-cavity surface-emitting laser: 850~nm class, MMF SR links.}

\declaretermsnip{LIV}{Light--current--voltage curve: threshold, slope efficiency, kink-free
  range, and thermal rollover versus bias.}

\declaretermsnip{SMSR}{Side-mode suppression ratio: power difference (dB) between the lasing
  mode and the strongest side mode (OSA).}

\declaretermsnip{RIN}{Relative intensity noise: laser amplitude noise spectral density
  (dB/Hz). Specs often quote stressed $\mathrm{RIN}_x\mathrm{OMA}$ under fixed ORL.}

\declaretermsnip{APC}{Automatic power control: feedback loop from a monitor PD that
  holds average optical power; slow loop BW keeps it out of the RIN band.}

\declaretermsnip{ELSFP}{External Laser Small Form-Factor Pluggable (OIF): faceplate CW laser
  module for CPO; 24-pin card edge, CMIS, blind-mate MT.}

\declaretermsnip{Arrhenius}{Life-acceleration model: temperature stress multiplies wear-out
  by $\exp[(E_a/k)(1/T_\mathrm{use}-1/T_\mathrm{stress})]$.}

\declaretermsnip{wavelengthLocker}{Wavelength locker: etalon + dual-PD error signal that
  trims laser TEC or current onto a grid slot (common on discrete DFB/EML).}

\declaretermsnip{thermalCrosstalk}{Thermal crosstalk: heater or self-heating on one ring
  shifts neighbors' resonances; worst case is full-bank traffic at max case $T$.}

\declaretermsnip{ACC}{Active copper cable: redriver (CTLE+VGA) inside cable; extends DAC reach.}

\declaretermsnip{AEC}{Active electrical cable: retimer (EQ+CDR) inside cable; longer twinax.}

\declaretermsnip{RLM}{Relative level mismatch: PAM4 level spacing uniformity. 1.0 is
  perfectly even; CEI often requires $\ge$0.95.}

\declaretermsnip{OMA}{Optical modulation amplitude: outer PAM4 level swing
  $P_1-P_0$.}

\declaretermsnip{COM}{Channel operating margin: statistical SNR of the fully equalized
  link. CEI go/no-go is typically COM $\ge$3~dB.}

\declaretermsnip{ISI}{Inter-symbol interference: energy from neighboring symbols
  overlapping the current unit interval.}

\declaretermsnip{DFE}{Decision-feedback equalizer: cancels post-cursor ISI using past
  symbol decisions.}

\declaretermsnip{CTLE}{Continuous-time linear equalizer: analog high-frequency boost in
  SerDes or TIA.}

\declaretermsnip{FFE}{Feed-forward equalizer: FIR filter. TDECQ applies a bounded FFE as
  the reference equalizer.}

\declaretermsnip{SNDR}{Signal-to-noise-and-distortion ratio: transmitter quality ceiling
  including nonlinearity.}

\declaretermsnip{UI}{Unit interval: one symbol period ($=1/R_{\mathrm{sym}}$). Jitter
  specs are in UI.}

\declaretermsnip{PMD}{Physical medium dependent: the optical transceiver (Tx/Rx) in
  IEEE 802.3.}

\declaretermsnip{PCS}{Physical coding sublayer: 64B/66B line coding, 256B/257B
  transcoding, scrambling, and RS-FEC. Maps MAC data to the coded line stream.}

\declaretermsnip{PMA}{Physical medium attachment: serialization, lane multiplexing,
  and clock recovery. The SerDes lives here.}

\declaretermsnip{AUI}{Attachment unit interface: the electrical serial lanes between
  host and module (e.g.\ 400GAUI-4, at the faceplate cage).}

\declaretermsnip{xMII}{Media-independent interface: a wide parallel bus between MAC and
  PCS inside the chip (e.g.\ 400GMII).}

\declaretermsnip{LPO}{Linear pluggable optics: no retimer/DSP in the module; host SerDes
  drives the modulator directly.}

\declaretermsnip{LRO}{Linear retimed optics: retimer in the module but no heavy DSP;
  lighter than fully retimed pluggables.}

\declaretermsnip{XSR}{CEI extra-short reach: die-to-die or die-to-OE, mm-scale traces.}

\declaretermsnip{CPC}{Co-packaged copper: copper I/O at the package edge (OIF map);
  feeds faceplate cages or short copper without an on-package optical engine.}

\declaretermsnip{VSR}{CEI very short reach: host PCB plus one module connector
  ($\sim$200~mm).}

\declaretermsnip{MR}{CEI medium reach: longer host plus module channel than VSR.}

\declaretermsnip{LR}{CEI long reach: longest CEI copper class before optics take over.}

\declaretermsnip{FR8DR8}{802.3 800G single-mode reaches: FR8 = 2~km, DR8 = 500~m
  (eight lanes per direction).}

\declaretermsnip{SerDes}{Serializer/deserializer: high-speed I/O that maps parallel
  data to PAM4 on the wire.}

\declaretermsnip{dspserdes}{DSP-based SerDes: ADC on the Rx lane, digital FFE/DFE and
  timing recovery, DAC on Tx. Standard above $\sim$50~GBd; the ``DSP'' that cancels ISI
  at 224G lives inside the SerDes.}

\declaretermsnip{GBd}{Gigabaud: $10^9$ symbols per second. Line rate = baud
  $\times$ bits per symbol.}

\declaretermsnip{PHY}{Physical layer: PCS plus PMD; maps MAC frames to the line code.}

\declaretermsnip{DCA}{Digital communication analyzer: sampling oscilloscope used for
  eye diagrams and TDECQ.}

\declaretermsnip{CDR}{Clock-data recovery: extracts bit clock from the data stream;
  present in retimed modules, absent in LPO.}

\declaretermsnip{EML}{Electro-absorption modulated laser: DFB plus EAM on one InP chip;
  dominant 100--200G/lane pluggable Tx.}

\declaretermsnip{UEC}{Ultra Ethernet Consortium: open Ethernet stack for AI/HPC scale-out (UET transport, PHY extensions).}
\declaretermsnip{UALink}{Ultra Accelerator Link: open scale-up fabric for accelerator load/store (200G/lane class).}
